\title{xLSTM: Extended Long Short-Term Memory}

\begin{abstract}
Long Short-Term Memory (LSTM) networks, pioneered in the 1990s through the constant error carousel and gating mechanisms, have proven fundamental to a host of major deep learning advances and notably formed the basis for the earliest Large Language Models (LLMs). Despite their historical importance, LSTMs were eventually surpassed by the emergence of Transformer architectures, which exploited scalable self-attention mechanisms to achieve superior performance in large-scale settings. This paper seeks to revisit LSTM architectures, posing the central question: What language modeling capabilities can be achieved when LSTMs are expanded to billions of parameters, incorporate recent advances inspired by LLM developments, and systematically address longstanding drawbacks? Our contributions are twofold. First, we propose exponential gating equipped with effective normalization and stabilization methods. Second, we revise LSTM memory components, introducing: (i) sLSTM, defined by a scalar-valued memory, a scalar update mechanism, and a novel memory mixing strategy, and (ii) mLSTM, a design that achieves full parallelization with a matrix-based memory and a covariance-driven update formulation. By embedding these new LSTM variants within residual-block frameworks, we obtain xLSTM blocks, which can be composed to build deep xLSTM architectures. The adoption of exponential gating alongside these redesigned memory structures substantially enhances xLSTM performance, allowing these models to compete with— and in some cases match or exceed— the capabilities of leading Transformer and State Space Model approaches in both performance metrics and scaling behavior.\\
Code available at: \url{https://github.com/NX-AI/xlstm}
\end{abstract}

\section{Introduction}
\label{sec:intro}

The concepts underpinning Long Short-Term Memory (LSTM)~\citep{Hochreiter:91,Hochreiter:97nips,Hochreiter:97}, such as the constant error carousel and the gating mechanisms, were developed as solutions to the vanishing gradient issue present in recurrent neural networks~\citep{Hochreiter:91,Hochreiter:00book}:

\begin{align}
\label{eq:lstm_idea}
  c_t \ = \ f_t \ c_{t-1} \ + \ 
          i_t \  z_t \ , \quad  
          h_t \ = \ o_t \ \psi({c_t}) \ . 
\end{align}

The constant error carousel functions by adding the cell inputs $z_t$ to the preceding cell state $c_{t-1}$ (depicted in green), with this addition being regulated by sigmoid gates (colored blue). Specifically, the forget gate $f_t$ and the input gate $i_t$ govern how this update is performed. The output gate $o_t$, on the other hand, dictates the cell's output, corresponding to the hidden state $h_t$. The value of the cell state is first normalized or squashed by $\psi$, after which the output gate determines the final hidden state.

LSTMs have achieved notable success across a broad spectrum of fields \citep{Hochreiter:01,Hochreiter:07,Schmidhuber:15} and were the dominant models for text generation prior to the emergence of Transformers in 2017~\citep{Vaswani:17}. Their efficacy has been established through a variety of sequence-oriented tasks, including but not limited to text generation~\citep{Graves:13,Karpathy:15blog}, handwriting synthesis~\citep{Graves:13}, sequence-to-sequence translation~\citep{Sutskever:14nips}, computer program analysis~\citep{Zaremba:14arxiv}, image caption generation~\citep{Karpathy:15,Hossain:19}, code generation~\citep{Karpathy:15blog}, rainfall-runoff forecasting~\citep{Kratzert:18,Kratzert:19}, and constructing hydrological models for flood prediction~\citep{Nearing:24}. Within the domain of reinforcement learning, LSTMs have established themselves as the most effective sequence architectures, as demonstrated by their integration in models such as AlphaStar for StarCraft II~\citep{Vinyals:17short}, OpenAI Five for Dota 2~\citep{OpenAI:19}, and the magnetic controller for nuclear fusion reactors~\citep{Degrave:22short}. Their particular strength lies in learning high-level abstractions—skilfully extracting and encoding semantic content in their cell states~\citep{Karpathy:15blog}—which is highlighted by the identification of specialized neurons such as number and syntax neurons~\citep{Lakretz:19}, linguistic neurons~\citep{Bau:19}, and sentiment neurons~\citep{Radford:17arxiv}. Notably, LSTMs continue to be deployed in critical real-world applications~\citep{Degrave:22short,Nearing:24} and continue to remain resilient and relevant over time.

Although LSTMs have achieved notable success, they exhibit three principal shortcomings: (i) Inability to amend prior storage decisions. To illustrate this constraint, we refer to the {\it Nearest Neighbor Search} task (see also Appendix~\ref{sec:appNearestNeighborSearch}). Given a reference vector, the process requires iterating through the sequence to identify and return the value associated with the most similar vector at the end of the sequence. The left side of Figure~\ref{fig:lstmProblems} depicts the mean squared error for this task. LSTM architectures encounter difficulty updating stored information when encountering a more similar vector later in the sequence. In contrast, our proposed xLSTM overcomes this with exponential gating mechanisms. (ii) Restricted capacity for storage, as information is forced into scalar-valued cell states. This issue is demonstrated through the {\it Rare Token Prediction} experiment: as shown in the right panel of Figure~\ref{fig:lstmProblems}, the token prediction perplexity for various token frequencies in Wikitext-103~\citep{Merity:17} reveals that LSTM performance deteriorates on infrequent tokens, attributable to these storage constraints. The xLSTM addresses this shortcoming with a matrix-based memory. (iii) Absence of parallelizability, caused by the entanglement of memory via hidden-to-hidden state dependencies across time steps, necessitating strictly sequential computation.

The shortcomings of LSTMs have contributed to the rise of Transformers~\citep{Vaswani:17} within the field of language modeling. An open question is what level of performance can be attained in language modeling tasks if these limitations are addressed and LSTMs are scaled up to match the capacities of today’s Large Language Models.

\section{Extended Long Short-Term Memory}
\label{sec:xLSTM}

In order to address the shortcomings inherent in LSTM architectures, the Extended Long Short-Term Memory (xLSTM) framework implements two principal alterations to the foundational concept described in Equation~\eqref{eq:lstm_idea}. Specifically, xLSTM incorporates exponential gating mechanisms as well as innovative memory configurations. These enhancements give rise to two distinct LSTM variants: (i) the sLSTM (see Section~\ref{sec:sLSTM}), which utilizes a scalar-valued memory component, relies on scalar updates, and applies memory mixing, and (ii) the mLSTM (see Section~\ref{sec:mLSTM}), which operates with matrix-valued memory and employs a covariance-based (outer product) update strategy that is amenable to full parallelization. Both variants augment the standard LSTM through the introduction of exponential gating. Unlike the sLSTM, the mLSTM forgoes memory mixing—effectively omitting hidden-to-hidden recurrent links—to facilitate parallel computation. Both sLSTM and mLSTM models can incorporate several memory cells; in the sLSTM, this allows for memory mixing between different cells. Furthermore, the sLSTM can support multiple heads, with memory mixing limited to interactions among cells within each head rather than across distinct heads. This extension, enabled by the combination of multiple heads and exponential gating, presents a novel approach to memory mixing in recurrent networks. In the case of mLSTM, the notions of multiple heads and multiple cells are functionally identical.

By embedding these novel LSTM variants within residual block modules, we form what are termed xLSTM blocks (refer to Section~\ref{sec:xLSTMblock}). Sequential residual stacking of these xLSTM blocks leads to the construction of xLSTM architectures (see Section~\ref{sec:xLSTMarchitecure}). Figure~\ref{fig:xlstm_sketch} illustrates the structure and components of the xLSTM architecture.

\subsection{Review of the Long Short-Term Memory}
\label{sec:orgLSTM}

The foundational concept behind LSTM networks~\citep{Hochreiter:91,Hochreiter:97nips,Hochreiter:97} features the scalar memory cell, which serves as the core element responsible for computation and memory, specifically designed to circumvent the vanishing gradient problem~\citep{Hochreiter:91,Hochreiter:00book} via the constant error carousel realized through updates to the cell state. The architecture equips each memory cell with three types of gates: the input, output, and the forget gate, the latter of which was proposed by \citet{Gers:00}. The following equations define how the LSTM memory cell is updated at each time step $t$:

\begin{align}
c_t \ &= \  f_t \ c_{t-1} \ + \ i_t \ z_t &  & &\text{cell state} \\
h_t  \ &= \ o_t \ \tilde{h}_t \ , & \tilde{h}_t \ &= \ \psi \left( c_t\right)
  &\text{hidden state} \\
z_t \ &= \ \varphi \left( \tilde{z}_t \right) \ , 
  &\tilde{z}_t \ &=  \ \mathbf{w}^\top_{z} \ \mathbf{x}_t \ + \
  r_{z}  h_{t-1} \ + \  b_{z} \ \
  &\text{cell input} \\
i_t \ &= \ \sigma \left( \tilde{i}_t  \right) \ , 
  &\tilde{i}_t \ &= \ \mathbf{w}^\top_{i} \ \mathbf{x}_t \ + \
  r_{i}  \ h_{t-1} \ + \  b_{i} \ \
  &\text{input gate} \\
f_t \ &= \ \sigma \left(  \tilde{f}_t \right) \ , 
  &\tilde{f}_t \ &= \ \mathbf{w}^\top_{f} \ \mathbf{x}_t  \ + \
  r_{f}  \ h_{t-1} \ + \  b_{f} \ \
  &\text{forget gate} \\
o_t \ &= \ \sigma \left( \tilde{o}_t \right) \ , 
  &\tilde{o}_t  \ &= \ \mathbf{w}^\top_{o} \ \mathbf{x}_t \ + \
  r_{o}  \ h_{t-1} \ + \  b_{o} \ \
  &\text{output gate} 
\end{align}

The vectors $\mathbf{w}_{z}$, $\mathbf{w}_{i}$, $\mathbf{w}_{f}$, and $\mathbf{w}_{o}$ represent the sets of weights connecting the inputs $\mathbf{x}_t$ to the cell input, input gate, forget gate, and output gate, respectively. The terms $r_{z}$, $r_{i}$, $r_{f}$, and $r_{o}$ denote the recurrent connections linking the previous hidden state $h_{t-1}$ to these same components within the cell. Corresponding biases are given by $b_{z}$, $b_{i}$, $b_{f}$, and $b_{o}$. Activation functions $\varphi$ and $\psi$, most commonly implemented as $\tanh$, serve for the cell input and the transformation of the hidden state. The activation $\psi$ is specifically applied to constrain the values of the cell state, which would otherwise grow without bound. Each gate in the model utilizes the sigmoid activation, defined as $\sigma \left( x \right)= 1/(1+\exp(-x))$. Subsequent extensions aggregate multiple scalar memory cells into the vector $\mathbf{c}_t \in \mathbb{R}^{d}$, permitting the use of recurrent matrices $\mathbf{R} \in \mathbb{R}^{d \times d}$ to integrate outputs from multiple memory cells~\citep{Greff:15}; see Appendix~\ref{sec:apporgLSTM} for additional explanation. Experiments involving ablation have highlighted that the memory cell’s functionality depends on the presence of all its components~\citep{Greff:15}.

\subsection{sLSTM}
\label{sec:sLSTM}

To enable LSTMs to reconsider storage choices, we propose the incorporation of exponential gates (red), along with techniques for normalization and stabilization. Specifically, the input and forget gates may utilize exponential activation functions. Regarding normalization, we employ a normalizer state designed to accumulate the sum of the input gate multiplied by the products of all subsequent forget gates.

The scalar sLSTM forward pass is:

\begin{align}
\label{eq:slstmforward}
c_t \ &= \  f_t \ c_{t-1} \ + \ i_t \ z_t & & 
 &\text{cell state} \\
n_t \ &= \  f_t \ n_{t-1} \ + \ i_t  & & 
 &\text{normalizer state} \\
h_t \ &= \ o_t \ \tilde{h}_t \ ,
  &\tilde{h}_t \ &= \ c_t / n_t
&\text{hidden state} \\
z_t \ &= \ \varphi \left( \tilde{z}_t \right) \ , 
  &\tilde{z}_t \ &=  \ \mathbf{w}^\top_{z} \ \mathbf{x}_t \ + \
  r_{z}  h_{t-1} \ + \  b_{z}
  &\text{cell input} \\
i_t \ &= \ \!\exp \left( \tilde{i}_t  \right) \ , 
  &\tilde{i}_t \ &= \ \mathbf{w}^\top_{i} \ \mathbf{x}_t \ + \
  r_{i}  \ h_{t-1} \ + \  b_{i}
  &\text{input gate} \\
f_t \ &= \ \sigma \left(  \tilde{f}_t \right) \ \text{OR} \
\!\exp \left(  \tilde{f}_t \right) \ ,
  &\tilde{f}_t \ &= \ \mathbf{w}^\top_{f} \ \mathbf{x}_t  \ + \
  r_{f}  \ h_{t-1} \ + \  b_{f}
  &\text{forget gate} \\
o_t \ &= \ \sigma \left( \tilde{o}_t \right) \ , 
  &\tilde{o}_t  \ &= \ \mathbf{w}^\top_{o} \ \mathbf{x}_t \ + \
  r_{o}  \ h_{t-1} \ + \  b_{o}
  &\text{output gate} 
\end{align}

The traditional gating mechanisms from LSTM models—including those reliant on input or hidden states along with a bias term—are incorporated into the revised architectures. Since exponential activations have the potential to generate excessively large outputs, resulting in numerical overflow, we incorporate an auxiliary state \mbox{$m_t$~\citep{Milakov:18arxiv}} to enhance the stability of the gates:

\begin{align}
\label{eq:slstmstabil}
m_t \ &= \ \max \left( \log ( f_t ) + m_{t-1} , \log ( i_t ) \right) &\text{stabilizer state} \\
i'_t \ &= \ \!\exp \left( \log \left( i_t \right) - m_t \right) \ = \ \!\exp \left( \tilde{i}_t - m_t \right) \ \, 
  &\text{stabil. input gate} \\
f'_t \ &= \ \!\exp \left( \log \left( f_t \right) + m_{t-1} - m_t \right) \ \, 
  &\text{stabil. forget gate}
\end{align}

In Appendix~\ref{sec:appsLSTM}, we demonstrate that substituting $f_t$ with $f'_t$ and $i_t$ with $i'_t$ during the forward pass has no effect on either the network's output or the loss gradients with respect to the parameters.

\paragraph{New Memory Mixing.} Similar to standard LSTM architectures, sLSTM supports the utilization of several memory cells (refer to Appendix~\ref{sec:appsLSTM} for additional details). This configuration allows for the blending of memory information through the recurrent matrices $\mathbf{R}_{\mathbf{z}}$, $\mathbf{R}_{\mathbf{i}}$, $\mathbf{R}_{\mathbf{f}}$, and $\mathbf{R}_{\mathbf{o}}$ which mediate the flow from the hidden state $\mathbf{h}$ to the memory cell input $\mathbf{z}$ and the gating units $\mathbf{i}$, $\mathbf{f}$, $\mathbf{o}$. What differentiates the updated memory mixing in this context is the incorporation of exponential gating. Each head in the revised sLSTM architecture can support memory mixing independently, though memory interactions do not occur between different heads. The combination of head-specific processing in sLSTM and exponential gating thus leads to a novel approach to integrating and updating memory.

\subsection{mLSTM}
\label{sec:mLSTM}

In order to augment the storage capabilities of LSTM networks, we replace the traditional memory cell, represented by a scalar $c \in \mathbb{R}$, with a matrix form $\mathbf{C} \in \mathbb{R}^{d \times d}$. Consequently, memory retrieval operations become matrix multiplications. At each time step $t$, we seek to encode a pair of vectors—the key $\mathbf{k}_t \in \mathbb{R}^d$ and the value $\mathbf{v}_t \in \mathbb{R}^d$—following the terminology established in the Transformer framework. Subsequently, at a future time $t + \tau$, a query vector $\mathbf{q}_{t+\tau} \in \mathbb{R}^d$ is utilized to access the value $\mathbf{v}_t$. This approach aligns with the paradigm of Bidirectional Associative Memories (BAMs) as described in previous works \citep{Kohonen:72,Anderson:72,Nakano:72,Anderson:77}.

The covariance update rule~\citep{Sejnowski:77,Dayan:91} for storing a key-value pair is

\begin{align}
\mathbf{C}_t \ &= \ \mathbf{C}_{t-1} \ + \ \mathbf{v}_{t} \ \mathbf{k}_t^\top \ .
\end{align}

We posit that inputs undergo layer normalization prior to being mapped onto key and value representations, resulting in these inputs possessing zero mean. According to~\citep{Dayan:91}, the covariance update mechanism achieves optimality in maximizing the separability among retrieved binary vectors—this is directly related to optimizing the signal-to-noise ratio. While greater separability could be obtained by restricting retrievals to pairwise associations, such a choice necessitates embracing the quadratic computational cost characteristic of attention mechanisms~\citep{Krotov:16,Krotov:17,Ramsauer:21}. Moreover, the covariance update method aligns with the principles underlying Fast Weight Programmers~\citep{Schmidhuber:92ncfastweights,Schlag:21}, for which subsequent work introduced fixed coefficients: a decay rate applied to $\mathbf{C}_{t-1}$, and a learning rate modulating 
$\mathbf{v}_{t} \mathbf{k}_t ^\top$~\citep{Ba:16}. Adopting this perspective, we embed the covariance update strategy within the LSTM architecture, drawing parallels between the forget gate and the decay rate, the input gate and the learning rate, and viewing the output gate as responsible for modulating the final retrieved vector.

In this matrix memory model, the normalizer state consists of a sum over the key vectors, each multiplied by both its associated input gate and the cumulative product of all subsequent forget gates. This configuration enables the normalizer to track the cumulative influence exerted by the various gates. Because the dot product between a query vector and the normalizer can approach zero, we instead consider its absolute value and apply a minimum threshold (commonly set to 1.0), consistent with earlier approaches~\citep{Sun:23arxiv}. The forward computation for mLSTM proceeds as follows:

\begin{align}
\label{eq:mlstm_recurrent_begin}
\mathbf{C}_t \ &= \  f_t \ \mathbf{C}_{t-1} \ + \ 
  i_t \ \mathbf{v}_t \ \mathbf{k}_t^\top & &  &\text{cell state} \\
\mathbf{n}_t \ &= \  f_t \ \mathbf{n}_{t-1} \ + \ 
  i_t \ \mathbf{k}_t & &  &\text{normalizer state} \\
\mathbf{h}_t  \ &= \ \mathbf{o}_t \ \odot \ \tilde{\mathbf{h}}_t
\ , \qquad \qquad 
\tilde{\mathbf{h}}_t \ = \   \mathbf{C}_t \mathbf{q}_t \ / \ 
  \max \left\{ |\mathbf{n}_t^\top \mathbf{q}_t|, 1 \right\} 
   &&&\text{hidden state} \\
\mathbf{q}_t \ &= \ \mathbf{W}_q \ \mathbf{x}_t \ + \ \mathbf{b}_q  & &  &\text{query input} \\
\mathbf{k}_t \ &= \ \frac{1}{\sqrt{d}} \mathbf{W}_k \ \mathbf{x}_t \ + \ \mathbf{b}_k  & &  &\text{key input} \\
\mathbf{v}_t \ &= \ \mathbf{W}_v \ \mathbf{x}_t \ + \ \mathbf{b}_v  & &  &\text{value input} \\
i_t \ &= \ \exp \! \left(\tilde{i}_t\right) \,, 
\qquad \qquad \qquad 
  \tilde{i}_t \ = \ \mathbf{w}^\top_{i} \ \mathbf{x}_t \ + \  b_{i} \ \ \ 
  &&&\text{input gate} \\
f_t \ &= \ \sigma \big( \tilde{f}_t \big) \ \text{OR} \ \exp \big( \tilde{f}_t \big), 
\qquad
\tilde{f}_t = \mathbf{w}^\top_{f} \mathbf{x}_t + b_{f}
  &&& \text{forget gate} \\
\label{eq:mlstm_recurrent_end}
\mathbf{o}_t \ &= \ \sigma \left( \tilde{\mathbf{o}}_t \right) \ , \qquad \qquad \qquad 
  \tilde{\mathbf{o}}_t  = \mathbf{W}_{\mathbf{o}} \mathbf{x}_t + 
  \mathbf{b}_{\mathbf{o}}
  &&&\text{output gate} 
\end{align}

mLSTM is capable of incorporating several memory cells, analogous to the architecture of standard LSTM. In the case of mLSTM, utilizing multiple heads and having multiple cells are functionally indistinguishable because there is no mixing of memory content. To ensure the exponential gates in mLSTM remain stable, we employ the identical stabilization methods applied to sLSTM (refer to Equation~\ref{eq:slstmstabil}). 
Given the absence of memory mixing within the mLSTM, its recurrence relations can be recast in a parallelized formulation. Additional information is available in Appendix~\ref{sec:appmLSTM}.

\subsection{xLSTM Architecture}
\label{sec:xLSTMarch}

\paragraph{xLSTM Blocks.}
\label{sec:xLSTMblock}
The xLSTM block is designed to capture and compress prior information through nonlinear transformations in a high-dimensional space, aiming to more effectively distinguish between different historical contexts. Distinguishing histories is essential for accurately forecasting the subsequent element in a sequence, such as the next token. Our approach builds upon Cover’s Theorem~\citep{Cover:65}, which posits that nonlinear mappings into higher-dimensional spaces tend to enhance the linear separability of data points compared to their separability in the original space. We investigate two forms of residual block designs: (i) a residual block employing post up-projection (inspired by Transformers), which initially forms a nonlinear representation of the historical context within the base space, then projects this summary linearly to a higher-dimensional space, activates it nonlinearly, and subsequently projects it back linearly to the initial dimension; refer to the left panel of Figure~\ref{fig:Backbones} and the third column of Figure~\ref{fig:xlstm_sketch}. A detailed schematic is provided in Figure~\ref{fig:appsLSTM_detailed} in the appendix. (ii) A residual block adopting pre up-projection (following the structure of State Space Models), which first linearly increases the dimensionality,

It captures historical information in a non-linear fashion within a high-dimensional space, subsequently performing a linear projection back into the original space. When an xLSTM block incorporates an sLSTM, the post up-projection block is utilized. Conversely, if the xLSTM block uses an mLSTM, the pre up-projection block is applied because the memory capacity increases when operating in the high-dimensional space. For further clarification, please consult the left panel of Figure~\ref{fig:Backbones}, the third column of Figure~\ref{fig:xlstm_sketch}, or Figure~\ref{fig:appsLSTM_detailed} in the appendix.

\paragraph{xLSTM Architecture. }
\label{sec:xLSTMarchitecure}
The xLSTM model is built by arranging its constituent blocks in a residual stack~\citep{Srivastava:15,He:16}. For this architecture, we adopt the widely prevalent pre-LayerNorm~\citep{Ba:16b} residual backbone architecture, following the standard approach in state-of-the-art Large Language Models. An illustration is provided in the final column of Figure~\ref{fig:xlstm_sketch}.

\subsection{Memory and Speed Considerations}

In contrast to Transformers, xLSTM networks maintain linear computational complexity and constant memory requirements as sequence length increases. Owing to the compressive nature of their memory, xLSTM models are particularly advantageous for industrial use cases and edge device deployment.

While the memory in mLSTM operates without additional parameters, it demands significant computational resources due to both the $d \times d$ memory matrix and its updates of the same dimensionality. Our approach balances the trade-off between maximizing memory capacity and managing computational demands. Importantly, because these operations can be efficiently parallelized on GPUs, the overall impact on real-world runtime is minimal.

Although mLSTM can be parallelized similarly to FlashAttention~\citep{Dao:22,Dao:23} and GLA~\citep{Yang:23arxiv}, sLSTM's inherent memory mixing (hidden-hidden connections) prevents such parallelization. Nonetheless, we have introduced an efficient CUDA implementation with optimizations extending to the GPU register level, resulting in sLSTM execution speeds generally less than twice as slow as those of mLSTM.

\section{Related Work}
\label{sec:related}

\paragraph{Linear Attention.} To address the Transformer’s quadratic computational cost with respect to context length, various techniques have been proposed to render attention mechanisms linear. For instance, the Synthesizer generates synthetic attention weights independently of explicit token-to-token comparisons~\citep{Tay:20arxiv}. The Linformer employs low-rank projection matrices to approximate self-attention, thereby achieving a linear reduction in computation~\citep{Wang:20arxiv}. Linear Transformer reformulates the attention computation so that it scales linearly with sequence length~\citep{Katharopoulos:20}. The Performer leverages positive orthogonal random features to efficiently approximate the softmax attention in a linear fashion~\citep{Choromanski:21}. Furthermore, alternatives such as Structured Global Convolution (SGConv)~\citep{Li:22} and the Hyena Hierarchy~\citep{Poli:23} replace traditional attention modules with highly efficient long convolutional operations.

\paragraph{State Space Models.} In recent times, State Space Models (SSMs) have garnered considerable attention due to their linear dependence on context length and competitive results relative to Transformers. Early advancements in this area include the Structured State Space sequence model (S4)~\citep{Gu:21}. Subsequent developments feature the Diagonal State Space (DSS) model~\citep{Gupta:22}, the Gated State Space (GSS) models~\citep{Mehta:22}, the S5 model~\citep{Smith:22}, the Bidirectional Gated SSM (BiGS)~\citep{Wang:22}, the H3 model~\citep{Fu:23}, and, more recently, Mamba~\citep{Gu:24arxiv}.

\paragraph{Recurrent Neural Networks.} Recent proposals have highlighted Recurrent Neural Networks (RNNs) as alternatives to Transformer architectures and attention mechanisms, owing to their linear computational complexity with respect to sequence length. Notably, RNNs incorporating Deep Linear Recurrent Units (LRUs) have achieved strong performance in language modeling tasks~\citep{Orvieto:23, De:24arxiv}. Hierarchically Gated Linear RNN (HGRN)~\citep{Qin:23} and its successor HGRN2~\citep{Qin:24arxiv} have also demonstrated efficacy. Another prominent example in the domain of large language models is the RWKV architecture~\citep{Peng:23arxivshort,Peng:24arxivshort}, which has yielded results on par with those of Transformer-based models.

\paragraph{Gating.}
Gating mechanisms, originally introduced in LSTM, have since been rediscovered and applied in numerous contemporary methods. Notable instances of gating include its integration into HGRN~\citep{Qin:23}, HGRN2~\citep{Qin:24arxiv}, Gated Linear Attention (GLA)~\citep{Yang:23arxiv}, Gated State Space (GSS) models~\citep{Mehta:22}, Bidirectional Gated SSM (BiGS)~\citep{Wang:22}, Moving Average Equipped Gated Attention (MEGA)~\citep{Ma:22}, RWKV~\citep{Peng:23arxivshort}, and Mamba~\citep{Gu:24arxiv}.

\paragraph{Covariance Update Rule.} In order to boost storage capacity, the mLSTM cell was augmented with a matrix memory governed by a covariance update rule. Comparable approaches that employ this kind of update procedure include Fast Weight Programmers \citep{Schmidhuber:92ncfastweights,Schlag:21}, RWKV-5 and RWKV-6~\citep{Peng:24arxivshort}, Retention~\citep{Sun:23arxiv}, Linear Transformer~\citep{Katharopoulos:20}, and HGRN2~\citep{Qin:24arxiv}.

\paragraph{Most Related.} In terms of conceptual similarities, xLSTM is most closely associated with models such as Retention~\citep{Sun:23arxiv}, RWKV~\citep{Peng:23arxivshort,Peng:24arxivshort}, and HGRN2~\citep{Qin:24arxiv}, all of which utilize ideas like matrix memory and gating. Unlike the newly introduced sLSTM, however, these prior models lack support for memory mixing. The ability to mix memories is crucial for addressing state tracking challenges, which makes LSTMs inherently more expressive than both State Space Models (SSMs) and Transformers~\citep{Merrill:24,Deletang:23}. Accurate state tracking is an important capability, for instance, when interpreting code or following entities across extended text passages.

\paragraph{Residually Stacking Architectures.} Similar to the majority of modern large-scale deep learning systems, the xLSTM models utilize residual stacking of modular blocks~\citep{Srivastava:15,He:16} in their design.

This stacking paradigm has been fundamental to the success of deep convolutional neural networks~\citep{He:16} and has also underpinned the architecture of Transformers~\citep{Vaswani:17}. Indeed, Transformers serve as the foundational technology behind Large Language Models (LLMs) such as GPT-3~\citep{Brown:20short}, ChatGPT~\citep{Schulman:22short}, GPT-4~\citep{Achiam:23gpt4short}, Megatron-LM~\citep{Shoeybi:19}, Gopher~\citep{Rae:21short}, ERNIE 3.0 Titan~\citep{Wang:21arxivshort}, GLaM~\citep{Du:21glamshort}, Chinese M6~\citep{Lin:21}, multilingual AlexaTM 20B~\citep{Soltan:22}, OPT~\citep{Zhang:22opt}, Chinchilla~\citep{Hoffmann:22short}, BLOOM~\citep{Scao:22arxivshort}, GLM-130B~\citep{Zeng:22arxivshort}, LaMDA~\citep{Thoppilan:22short}, PaLM~\citep{Chowdhery:22short}, Llama~\citep{Touvron:23arxiv}, and Gemini~\citep{GeminiTeam:23arxiv,Reid:24arxiv}.

\section{Experiments}
\label{sec:Exp}

This section presents an empirical analysis of xLSTM, placing particular emphasis on its performance in language modeling and benchmarking it against prevailing approaches. Section~\ref{sec:ExpSyntethic} is dedicated to examining the unique strengths of xLSTM through a suite of synthetic tasks. Section~\ref{sec:ExpComparison} details a comparative study of the validation perplexity across a range of contemporary language modeling techniques, all of which are trained on a corpus of 15B tokens from SlimPajama~\citep{Soboleva:23}. This section further includes ablation experiments specific to xLSTM. Additionally, we explore how these models scale by conducting analyses similar to those found in \citet{Kaplan:20} and \citet{Brown:20short}.

Section~\ref{sec:ExpLanguage} provides an extended investigation into language modeling, focusing on xLSTM and the strongest baselines as identified in Section~\ref{sec:ExpComparison}, after training them on an expanded 300B token subset of SlimPajama~\citep{Soboleva:23}. This evaluation proceeds along four axes: (1) an examination of the methods’ ability to generalize to longer context lengths; (2) standard validation perplexity metrics and performance assessment on downstream tasks~\citep{Sutawika:24}; (3) an evaluation involving 571 distinct text domains drawn from the PALOMA language benchmark dataset~\citep{Magnusson:23arxivshort}; and (4) a renewed analysis of model scaling behaviors, this time leveraging a training dataset that is twenty times larger.

\label{sec:xlstm_blocks_notation}
Throughout all our experimental setups, we denote the ratio of mLSTM-based to sLSTM-based xLSTM blocks as xLSTM[$a$:$b$], where $a/b$ indicates the proportion between the two types. As an illustration, xLSTM[7:1] specifies a configuration in which seven out of eight blocks are built on mLSTM, with one block utilizing sLSTM. When considering the commonly used total of 48 blocks, this would comprise 42 mLSTM-based and 6 sLSTM-based blocks. In addition, each experiment consistently employs a pre up-projection block with mLSTM modules and a post up-projection block with sLSTM modules.

\subsection{Synthetic Tasks and Long Range Arena}
\label{sec:ExpSyntethic}
We begin by evaluating xLSTM’s novel exponential gating mechanism coupled with memory mixing on a suite of formal language benchmarks~\citep{Deletang:23}. Subsequently, we investigate how the newly introduced matrix memory in xLSTM performs in the context of the Multi-Query Associative Recall task~\citep{Arora:23arxiv}. Lastly, we examine xLSTM’s capability to manage long sequence data using the Long Range Arena framework~\citep{Tay:21}.

\paragraph{Assessment of xLSTM’s Exponential Gating with Memory Mixing.} We evaluate the newly introduced exponential gating mechanism with memory mixing in xLSTM, hypothesized to enhance its capability for state tracking tasks~\citep{Merrill:24, Merrill:23}. The formal language tasks from \citet{Deletang:23} are reimplemented and expanded to accommodate multi-length training, facilitating experiments in length extrapolation. For comprehensive task explanations and expanded performance metrics, refer to Appendix~\ref{sec:appExpSynthetic-formalLang}. 

Our analysis benchmarks xLSTM against various alternative architectures, including Transformers, State Space Models, and traditional Recurrent Neural Networks. Method accuracy is computed specifically for tokens central to task success and reported on a normalized scale from 0 (random) to 1 (optimal). We examine 2-block variants of several architectures, including: xLSTM[0:1] (pure sLSTM), xLSTM[1:0] (pure mLSTM), xLSTM[1:1], Llama, Mamba, RWKV, Retention, Hyena, LSTM, and a model using LSTM within Transformer blocks (LSTM (Block)). Figure~\ref{fig:formal-main} displays the empirical findings. Models such as Transformers and State Space Models lacking memory mixing—hence, incapable of effective state tracking—are unable to solve certain formal language problems, such as the parity task, which relies on regular grammar recognition. This outcome corroborates prior evidence indicating inherent limitations of Transformers and State Space Models relative to RNNs~\citep{Merrill:24, Merrill:23, Deletang:23}.

\paragraph{Test of xLSTM's Memory Capacities on Associative Recall Tasks.} This study assesses the matrix memory capabilities of xLSTM by measuring its memory capacity on the Multi-Query Associative Recall task~\citep{Arora:23arxiv}. In this benchmark, sequences are formed using randomly selected key-value pairs from a large vocabulary, requiring the model to memorize them for subsequent retrieval. To introduce greater challenge beyond the baseline task, we raise the number of key-value pairs to as many as 256 and extend the context window to a length of 2048. These adjustments enable a more comprehensive evaluation of memory constraints across model architectures. We conduct comparisons across 2-block variants of Llama, Mamba, RWKV-5, RWKV-6, xLSTM[1:1], and xLSTM[1:0]. Model performance is determined by the accuracy with which the stored pairs are recalled. Given that Transformer-based models (such as Llama) possess memory scaling exponentially with the coding dimension~\citep{Ramsauer:21}, they are taken as the benchmark for this task. The outcomes, summarized in Figure~\ref{fig:mqar-main}, indicate that xLSTM[1:1] surpasses other non-Transformer models, even at smaller model scales. Notably, the integration of the sLSTM block does not restrict, but rather augments, overall memory capacity, a finding especially marked in the most demanding scenario involving 256 key-value pairs. Further analyses and results can be found in Appendix~\ref{sec:appExpSynthetic-mqar}, where extrapolation experiments suggest that the superior memory traits of xLSTM facilitate generalization to contexts exceeding those encountered during training.

\paragraph{Test of xLSTM's Long Context Capabilities on Long Range Arena.} To evaluate the ability of xLSTM to manage extended sequences and broad context lengths, we benchmark it against various approaches using the Long Range Arena~\citep{Tay:21}. Across all benchmarks, xLSTM reliably attains strong results, indicating that the architecture is particularly adept at addressing the challenges associated with long-context tasks.

For more details, see Appendix~\ref{sec:appExpSynthetic-lra}.

\subsection{Method Comparison and Ablation Study}
\label{sec:ExpComparison}

In order to explore the primary research question of this study—namely, the capabilities of our proposed LSTM variants when applied at scale for language modeling—we conduct experiments where xLSTMs, Transformers, State Space Models, and additional techniques are each trained on 15B tokens from the SlimPajama dataset using an identical auto-regressive framework. The resulting models are evaluated on a shared validation set, and we further carry out ablation analyses focused on the xLSTM architectures.

\paragraph{Comparing xLSTM to Other Methods.} For this comparison, all models are trained on a dataset containing 15B tokens sourced from SlimPajama~\citep{Soboleva:23}. Model quality is assessed via perplexity on a dedicated validation set. The evaluation covers a range of approaches: the newly proposed xLSTM, GPT-3 (a Transformer-based model)~\citep{Brown:20short}, Llama (Transformer)~\citep{Touvron:23arxiv}, H3 (a type of State Space Model, SSM)~\citep{Fu:23}, Mamba (SSM)~\citep{Gu:24arxiv}, several iterations of RWKV (RNN variants: RWKV-4, RWKV-5, RWKV-6)~\citep{Peng:23arxivshort,Peng:24arxivshort}, GLA (linear Transformer)~\citep{Yang:23arxiv}, HGRN and HGRN2 (both RNN-based)~\citep{Qin:23,Qin:24arxiv}, along with RetNet (linear Transformer)~\citep{Sun:23arxiv}, Hyena (linear Transformer)~\citep{Poli:23}, and two xLSTM variants (xLSTM[1:0] and xLSTM[7:1], see Section~\ref{sec:xlstm_blocks_notation}). 

All models are trained using mixed precision; however, for RWKV-5, RWKV-6, GLA, and HGRN2, the implementation of mixed precision diverged from PyTorch's automated mixed precision framework (additional details are provided in Appendix Section~\ref{sec:appGenTrainProc}). The evaluation organizes these models into (a) Transformer architectures, (b) State Space Models (SSMs), and (c) a combined group of Recurrent Neural Networks (RNNs) and linear Transformer variants, with the latter replacing standard attention with linearized attention operations. Every architecture is standardized to match GPT-3's configuration, featuring approximately 350M parameters (embedding dimension 1024, 24 residual blocks), though only GPT-3 applies parameter sharing between the token and output embeddings—resulting in a reduced parameter count. According to results shown in Table~\ref{tab:spaj15b_model_results}, xLSTM yields the lowest validation perplexity, outperforming all other baselines. Further methodological and experimental information can be found in Appendix~\ref{sec:appExpComparison}. Scaling results are visualized in Figure~\ref{fig:temp_sclaw15B}, suggesting that xLSTM’s performance advantage is expected to persist with increased model capacity.

\paragraph{Ablation Studies.}
As shown in Table~\ref{tab:spaj15b_model_results} and Figure~\ref{fig:temp_sclaw15B}, xLSTM delivers outstanding performance in language modeling tasks with training on 15B tokens from SlimPajama. In order to isolate the impact of architectural modifications from LSTM to xLSTM, we incrementally transform a standard LSTM architecture into the xLSTM configuration. The process begins by embedding LSTM layers within pre-LayerNorm residual frameworks. In the next step, we introduce a post up-projection component. The final transformation stage incorporates exponential gating and matrix memory mechanisms.

The results are shown in Table~\ref{tab:ablstudies} (top). 

The ablation experiments indicate that both the exponential gating and the matrix memory substantially contribute to the observed gains in performance. Furthermore, given the centrality of gating in RNNs and State Space Models, we investigate alternative gating strategies. Results in Table~\ref{tab:ablstudies} (bottom) reveal that allowing each gate to be input-dependent and trainable consistently yields further improvements. Further analyses examining separate backbone elements can be found in Appendix~\ref{sec:appAblDetails}.

\subsection{xLSTM as Large Language Model}
\label{sec:ExpLanguage}

This research concludes with extensive language modeling experiments to evaluate xLSTM as a large language model (LLM). To this end, we scale up the training dataset and utilize 300B tokens sourced from SlimPajama. Comparable token counts are employed by other models such as Mamba~\citep{Gu:24arxiv} and Griffin~\citep{De:24arxiv}.

We benchmark xLSTM against RWKV-4, Llama, and Mamba, each representing a distinct methodological approach described in Section~\ref{sec:ExpComparison}. RWKV-4 was chosen as the RNN comparator because suitable training precision configurations for RWKV-5, RWKV-6, and HGRN2 were only established after initiating the 300B token experiments (further details are provided in Appendix~\ref{sec:appGenTrainProc}). Our experimental setup involves training multiple model sizes (125M, 350M, 760M, and 1.3B parameters), evaluating every model's ability to extrapolate to longer sequences, and measuring validation set performance. Additionally, we examine each model's results on downstream tasks, their language modeling proficiency across 571 text domains included in the PALOMA benchmark, and explore their scaling law characteristics.

\paragraph{Sequence Length Extrapolation.}
\label{sec:spaj300B_ctx_extrapolation}
We begin by evaluating the sequence length extrapolation abilities of large-scale, 1.3B parameter models, specifically xLSTM, RWKV-4, Llama, and Mamba. Each model is trained using a context length of 2048 and subsequently assessed on much longer contexts, reaching up to 16384 tokens. The performance outcomes are illustrated in Figure~\ref{fig:spaj300B_seq_extrapolation}. Unlike the other architectures, xLSTM consistently demonstrates low perplexity scores as context length increases.

\paragraph{Validation Perplexity and Downstream Tasks.}
\label{sec:spaj_downstream_eval}
Additionally, across every model scale, we assess the xLSTM, RWKV-4, Llama, and Mamba architectures on the SlimPajama validation set with respect to next token prediction. These models are also tested on downstream benchmarks designed to evaluate their common sense reasoning abilities.

The third column in Table~\ref{tab:lmeval} presents the perplexity scores on the validation set for various approaches. Across all model sizes, xLSTM[1:0] and xLSTM[7:1] consistently achieve the lowest perplexities, indicating superior validation set performance. The remaining columns in Table~\ref{tab:lmeval} summarize results on downstream tasks. xLSTM demonstrates leading performance on the majority of tasks and across all model scales, with Mamba surpassing it only occasionally on the ARC benchmark. Further information can be found in Appendix~\ref{sec:appExpLanguage}.

\paragraph{Performance on PALOMA Language Tasks.} Next, we evaluate models of various sizes, including xLSTM, RWKV-4, Llama, and Mamba, by assessing their next-token prediction capabilities on PALOMA language tasks~\citep{Magnusson:23arxivshort}. Performance is quantified using perplexity across 571 distinct text domains, spanning sources such as nytimes.com and the r/depression subreddit. Table~\ref{tab:ppleval} presents token prediction perplexities, with the initial seven columns corresponding to general language modeling benchmarks, and the final five capturing performance on more specific domain tests. On these language tasks, xLSTM[1:0] consistently surpasses xLSTM[7:1].

xLSTM[1:0] achieves a lower perplexity compared to Mamba in 568 of 571 text domains (99.5\%), outperforms Llama in 486 of 571 domains (85.1\%), and surpasses RWKV-4 in 570 of 571 cases (99.8\%). Further information is provided in Appendix~\ref{sec:appExpLanguage}.

\paragraph{Scaling Laws.} Fourth, we investigate the presence of power-law scaling, which is instrumental for predicting how performance evolves as model size increases~\citep{Kaplan:20,Brown:20short}. In Figure~\ref{fig:temp_sclaw300B}, the observed scaling patterns are illustrated. While all the models demonstrate analogous scaling tendencies, their performance curves are shifted relative to each other. RWKV-4 exhibits the lowest performance, trailed by Llama and subsequently Mamba. xLSTM surpasses Mamba by a margin comparable to the gap observed between Mamba and Llama. This scaling analysis suggests that, as models grow larger, xLSTM is likely to maintain an advantage over both Transformer-based and State-Space models.

\textbf{Generation Times and Maximal Throughput.} In Figures~\ref{fig:inference_time_generation} and~\ref{fig:inference_time_decoding_throughput} (left), we present results on generation latency and peak throughput, respectively, for our xLSTM models at the 1.3B parameter scale. These results are compared with HuggingFace implementations of models of comparable size, specifically Mamba, Llama, and RWKV, with the Llama baseline also incorporating a static key-value cache. At the point of conducting these experiments, neither full cache compilation for the Transformer-based model nor using \texttt{torch.compile} to compile the Mamba model was functional. All evaluated models were assessed under generation with a batch size of 1 and a pre-fill of 16 tokens, an experimental setting chosen to optimally benefit the Transformer baseline.

Figure~\ref{fig:inference_time_generation} illustrates that both xLSTM and other recurrent alternatives such as Mamba and RWKV-4 exhibit linear time complexity with respect to generation, whereas Llama displays quadratic scaling. For throughput assessment, we experimented with a range of batch sizes and prefill options specifically for the Llama model.

As shown in Figure~\ref{fig:inference_time_decoding_throughput} (right), xLSTM is capable of effectively utilizing much larger batch sizes than Llama, attributable to its constant memory footprint. This property enables xLSTM to achieve the best throughput among the evaluated architectures.

\section{Limitations}

(i) Unlike mLSTM, the memory mixing mechanism in sLSTM precludes the use of parallelizable computations, making rapid parallel execution unattainable. Despite this limitation, we engineered an efficient CUDA kernel for sLSTM that currently operates at less than twice the runtime of our parallel mLSTM kernel. (ii) Our mLSTM CUDA kernels remain unoptimized, resulting in performance approximately four times slower than FlashAttention or the scan operator utilized by Mamba. Future improvements akin to those seen with FlashAttention could yield significantly faster kernels. (iii) The matrix memory mechanism in mLSTM entails a substantial computational burden, since it requires manipulating $d \times d$ matrices. Nonetheless, because both memory updates and retrieval steps are parameter-free and can leverage parallel standard matrix operations, the real-time computational penalty due to the matrix memory remains minimal. (iv) Proper initialization of the forget gates is critical. (v) Since mLSTM's matrix memory is not tied to sequence length, longer input sequences could exceed the memory’s capacity as context grows. However, this does not seem to impede performance for context lengths up to 16k, as detailed in Section~\ref{sec:spaj300B_ctx_extrapolation}. (vi) Owing to the significant computational resources demanded for large language model experiments, we have not comprehensively optimized either the model architecture or the hyperparameters—particularly for more sizeable xLSTM variants. We expect that realizing the full capabilities of xLSTM will require a much more exhaustive tuning and architecture search effort.

\section{Conclusion}
Our investigation provides a partial response to the foundational question: What are the capabilities of language modeling when LSTM architectures are scaled to the order of billions of parameters? At present, the evidence demonstrates that such scaling allows LSTM to reach, and even match, the performance of leading approaches such as Transformers and State Space Models. Through the introduction of the xLSTM—an LSTM variant employing exponential gating with memory mixing and an innovatively restructured memory—we observe that xLSTM achieves results that are competitive with the most advanced techniques in language modeling. Analysis of scaling laws suggests that as xLSTM architectures are further expanded, they are poised to rival established Large Language Models based on Transformer designs. Beyond language modeling, xLSTM is well-positioned to make substantial contributions across domains including Reinforcement Learning, Time Series Forecasting, and physical system modeling.

\end{document}